\documentclass[Total.tex]{subfiles}

\begin{document}
	
	\chapter{Riemannian Metrics}
		\section{Definitions}
			\begin{Def}
				\begin{itemize}
					\item Let $M$ be a smooth manifold. A \textbf{Riemannian metric} on $M$ is a $2$-tensor field $g\in \mathcal{T}^2M$, whose value
					\begin{gather*}
						g_p:T_p M\times T_p M\rightarrow \RR
					\end{gather*}
					which is positive definite in the sense of $g_p(v,v)\geq 0$. 
					\item A \textbf{Riemannian manifold} is a pair $(M,g)$.
				\end{itemize}
			\end{Def}
			
			\textbf{Remark} Riemannian metric vs Metric Space
			
			\begin{example}
				(Eucilidean Metric) 
			\end{example}
			\subsubsection{Isometires}
			\subsubsection{Local Representation}
				\begin{Def}
					(Local Representation of Riemannian Metric)
				\end{Def}
				
				
				\begin{proposition}
					(Existence of Orthonormal Frame)
				\end{proposition}
				\begin{proof}
					内容...
				\end{proof}
				\begin{Def}
					(Unit Tangent Bundle)   Define the \textbf{unit tangent bundles}
					\begin{gather*}
						UTM:=\{(p,v)||v|_g=1\}
					\end{gather*}
				\end{Def}
				
				\begin{proposition}
					If $(M,g)$ is a Riemannian-Manifold with or without boundary, it unit tangent bundle $UTM$ is a \underline{smooth, properly embedded codimensional-1 submanifold with boundary in} $TM$, with
					\begin{gather*}
						\pi:UTM\rightarrow M\qquad \partial(UTM)=\pi^{-1}(\partial M)
					\end{gather*}
					Then, $UTM$ is connected$\Leftrightarrow$ $M$ is connected. $UTM$ is compact$\Leftrightarrow M$ is compact. 
				\end{proposition}
				\begin{proof}
					内容...
				\end{proof}
			\subsection{Riemannian Submanifold}
				\begin{lemma}
					Suppose $(\tilde{M},\tilde{g})$ is a Riemannian manifold with or without boundary. $M$ is a smooth manifold with or without boundary. And 
					\begin{gather*}
						F:M\rightarrow \tilde{M}
					\end{gather*}
					is a smooth map. The smooth 2-tensor field(pullback)
					\begin{gather*}
						g=F^*\tilde{g}
					\end{gather*}
					is a Riemannian metric iff $F$ is an immersion.
				\end{lemma}
				\begin{proof}
					内容...
				\end{proof}
		\section{Models of Riemannian Manifolds}
			\subsection{Symmetries}
			\subsection{Covering Space of Riemannian Manifold}	 
				
				\subsubsection{Eucilidean Space}
				\subsubsection{Sphere}
				\subsubsection{Hyperbolic Space}
				\subsubsection{Killing-Hopf Theorem}
					\begin{theorem}
						Suppose $(M,g)$ be compelete, simply connected Riemannian $n$-manifold with constant sectional curvature, $n\geq 2$.
						
						Then, $M$ is isometric to one of the model spaces $\RR^n,\mathbb{S}^n(R),\mathbb{H}^n(R)$.
					\end{theorem}
					\begin{proof}
						内容...
					\end{proof}
					
					\begin{corollary}
						内容...
					\end{corollary}
					\begin{proof}
						内容...
					\end{proof}
				\subsubsection{Space Form}
	\chapter{Riemannian Connections}
		\section{Affine Connections}
			\begin{Def}
				Define the \textbf{affine connections}
				\begin{gather*}
					\nabla:\mathfrak{X}(M)\times \mathfrak{X}(M)\rightarrow \mathfrak{X}(M)\\
					(X,Y)\mapsto 
				\end{gather*}
				With following properties:
				\begin{itemize}
					\item ($C^\infty(M)$-Linear) $\triangledown_{fX+gY}Z=f\triangledown_X Z+g\triangledown_{Y}Z$;
					\item($\RR$-Linear) $\triangledown_{X}(aY+bZ)=a\triangledown_X(Y)+b\triangledown_X(Z)$
					\item (Lebnitz)$\triangledown_{X}(fY)=f\triangledown_X Y+X(f)Y$
				\end{itemize}
				where $X,Y,Z\in \mathfrak{X}(M),f,g\in C^\infty(M)$.
			\end{Def}
			
			\begin{proposition}
				Let $M$ be a differentiable manifold with an affine connection $\triangledown$. There exists a \underline{unique} correspondnece which associates to a vector field $V$ along the differentiable curve
				\begin{gather*}
					c:I\rightarrow M
				\end{gather*}
				another vector field aloing $c$, called the \textbf{covariant derivative} of $V$ along $c$, such that:
				\begin{itemize}
					\item 
					\item 
					\item 
				\end{itemize}
			\end{proposition}
			\begin{proof}
				内容...
			\end{proof}
			\subsection{Parallel Vector Field}
		\section{Riemannian Metrics, Levi-Civita Theorems}
	\chapter{Geodesics}
		\section{The Geodesics Flow}
		\section{Minimalizing Properties}
	\chapter{Curvatures}
		\section{Curvatures}
		\section{Sectional Curvature}
		\section{Ricci Curvature and Scalar}
		\section{Tensors on Riemannian Manifolds}
		\section{Manifolds of Nonpositive Curvature}
			\subsection{Cartan-Hadamard Manifold}
			\begin{theorem}
				(Cartan-Hadamard)
			\end{theorem}
			\begin{proof}
				内容...
			\end{proof}
			\begin{Def}
				内容...
			\end{Def}
			\begin{proposition}
				内容...
			\end{proposition}
			\begin{proof}
				内容...
			\end{proof}
			
			\begin{proposition}
				内容...
			\end{proposition}
			\begin{proof}
				内容...
			\end{proof}
			\subsection{Cartan's Fixed Point Theorem}
			\subsection{Cartan's Torsion Theorem}
	\chapter{Appendix: Curvatures by Curves}
		\section{Classification of Curves: A Review}
			\begin{Def}
				Given a $C^1$-curve $c:[a,b]\rightarrow \RR^n$:
				
				\begin{itemize}
					\item $c$ is called \textbf{regular} if
					\begin{gather*}
						\Vert \dot{c}(t)\Vert\not=0,t\in [a,b]
					\end{gather*}
					\item $c$ is called \textbf{path-parametrized} if $\Vert \dot{c}(t)\Vert=1,t\in [a,b]$;
					\item $c$ is called \textbf{closed}, if $c(a)=c(b)$;
					\item $c$ is called \textbf{differentiable closed} if $c(a)=c(b);c'(a)=c'(b)$;
				\end{itemize}
			\end{Def}
			
			\begin{proposition}
				Suppose $c:[a,b]\rightarrow \RR^n$ be $C^k$-regular, then there exists a $C^k$-diffeomorphism
				\begin{gather*}
					\operatorname{Diff}^k([0,1],[a,b])\ni \varphi:[0,l]\rightarrow [a,b]\qquad \tilde{c}=c\circ \varphi
				\end{gather*}
				
				where $\tilde{c}$ is path-parametrized, $\Vert \dot{\tilde{c}}\Vert=1$;
			\end{proposition}
			\begin{proof}
				内容...
			\end{proof}
			
			\begin{Def}
				\begin{itemize}
					\item Given a regular $C^1$-curve $c:I\rightarrow \RR^n$, then
					\begin{gather*}
						T(t):=\frac{\dot{c}(t)}{\Vert \dot{c}(t)\Vert}
					\end{gather*}
					
					is called the unit tangential vector field. 
					\item Given a regular $C^2$-curve $c$. Then, define
					\begin{gather*}
						\kappa(t):=\frac{\Vert T'(t)\Vert}{\Vert \dot{c}(t)\Vert}
					\end{gather*}
					
					is called the \textbf{curvature} of $c$ at $t$;
				\end{itemize}
			\end{Def}
			
			As a matter of fact, $T_t c:\RR\cong T_t I\rightarrow T_{c(t)} \RR^n\cong \RR^n$, this is a linear map, then
			\begin{gather*}
				T(t)=T_tc(1)
			\end{gather*}
			
			\begin{proposition}
				Suppose $c$ be regular. Then,
				\begin{gather*}
					\kappa(t)=\frac{\Vert (\ddot{c}(t))^\perp \Vert}{\Vert \dot{c}(t)\Vert^2}
				\end{gather*}
				
				where $ (\ddot{c}(t))^\perp$ is the orthogonal projection onto hyperplane $T(t)^\perp$, i.e. $=\ddot{c}(t)-T(t)<T(t),\ddot{c}(t)>$;
			\end{proposition}
			\begin{proof}
				内容...
			\end{proof}
			
			\begin{corollary}
				Suppose $c$ be path-parametrized, then $\kappa(t)=\Vert \ddot{c}(t)\Vert$.
			\end{corollary}
			\begin{proof}
				
			\end{proof}
			
			\begin{proposition}
				(Curvature is not related to Parametrimization)
			\end{proposition}
			\begin{proof}
				内容...
			\end{proof}
			
			\begin{Def}
				Given a regular $C^2$ curve, with $\kappa(t_0)\not=0,t_0\in I$:
				\begin{itemize}
					\item 
					\item 
					\item 
				\end{itemize}
			\end{Def}
			
			
		\section{Curvature of Curve}
			\subsection{Total Curvature}
				\begin{lemma}
					内容...
				\end{lemma}
				\begin{proof}
					内容...
				\end{proof}
			\subsection{Curvature on Plane}
				\begin{proposition}
					We could calculating the unit tangent vector field and normal vector field as following
					\begin{Def}
						Define the spin operation 
						\begin{gather*}
							\begin{pmatrix}
								cos\theta&-sin\theta\\
								sin\theta&cos\theta
							\end{pmatrix}\qquad J:=\begin{pmatrix}
								0&-1\\
								1&0
							\end{pmatrix},\theta=\frac{\pi}{2}
						\end{gather*}
					\end{Def}
					
					Then, 
					\begin{gather*}
						N(t)=JT(t)
					\end{gather*}
				\end{proposition}
				\begin{proof}
				\end{proof}
				
				And the curvature on plane is 'orientable', because of the property of $\RR^2$: 
			\subsection{Knotted Curves}
			
		\section{Isoperimetric inequality}
			\begin{lemma}
				Suppose $c:I\rightarrow \RR^2$ be: simply closed, positive oriented, regular $C^1$-curve, and with coordinate
				\begin{gather*}
					c(t)=(x(t),y(t))
				\end{gather*}
				
				Then, for the area described by the curve, 
				\begin{gather*}
					A=-\int_0^l x'(t)y(t)=\int_0^l x(t)y'(t)dt
				\end{gather*}
			\end{lemma}
			\begin{proof}
				内容...
			\end{proof}
			
			\begin{proposition}
				(Isometric Inequality)
			\end{proposition}
	\chapter{Riemannian Submanifold}
	\chapter{Gauss-Bonnet Theorem}
	\chapter{Jacobi Field}
	\chapter{Comparison Theory}
	\chapter{Curvature and Topology}
	\chapter{Isomorphic Immersions}
	\chapter{Compelete Manifold, Hopf-Rinow and Hadamard Theorems}
	\chapter{Spaces of Constant Curvature}
	\chapter{Variations of Energy}
	\chapter{The Morse Index Theorem}
	\chapter{The Fundamental Group of Manifolds of Negative Curvatures}
	\chapter{The Sphere Theorem}
	\chapter{The Laplacian}
		This will be given in geometric analysis.
\end{document}